(1) 方程与无通量边界为
\[
c_t+u(r)c_z=D\{c_{zz}+r^{-1}(rc_r)_r\},\quad c_r(a,z,t)=0,\quad c_r(0,z,t)=0.
\]
定义 \(\langle q\rangle=2a^{-2}\int_0^a q(r)r\,dr,\ C=\langle c\rangle\)。径向扩散项的截面平均由壁面通量给出零，因此 \(C_t+\partial_z\langle uc\rangle=D C_{zz}\)；在轴向无流出边界时总示踪剂质量严格守恒。
(2) 首阶宏观方程 \(C_t=-U C_z+O(C_{zz})\)，代入渐近式比较 \(C_z\) 得
\[
D\,r^{-1}(r\chi')'=u-U=U(1-2r^2/a^2),\quad \chi'(0)=\chi'(a)=0.
\]
积分得到 \(\chi'=Ur(1-r^2/a^2)/(2D)\)，再积分得 \(\chi=(U/D)(r^2/4-r^4/(8a^2)+C_0)\)。用 \(\langle r^2\rangle=a^2/2,\ \langle r^4\rangle=a^4/3\) 与 \(\langle\chi\rangle=0\)，确定 \(C_0=-a^2/12\)；任意常数自由度由此唯一固定。
(3) 轴向总通量为 \(J=\langle uc\rangle-D C_z=UC+[\langle(u-U)\chi\rangle-D]C_z+\cdots\)，故 \(D_{\rm eff}=D-\langle(u-U)\chi\rangle\)。将胞元方程乘 \(\chi\) 并分部积分，边界项消失，得到
\[
-\langle(u-U)\chi\rangle=D\langle(\chi')^2\rangle
=\frac{U^2}{4D}\left\langle r^2(1-r^2/a^2)^2\right\rangle
=\frac{U^2a^2}{48D}.
\]
其中径向矩 \(\langle r^{2j}\rangle=a^{2j}/(j+1)\)，故增强项确为非负，且系数可独立复算。
(4) 宏观长时方程 \(C_t+UC_z=D_{\rm eff}C_{zz}\) 的质量为 \(M_0\) 的脉冲解为
\[
C(z,t)\simeq\frac{M_0}{\sqrt{4\pi D_{\rm eff}t}}\exp[-(z-Ut)^2/(4D_{\rm eff}t)].
\]
质心 \(\bar z=Ut\)，方差 \({\rm Var}(z)=2D_{\rm eff}t\)。当 \(U\to0\) 或 \(a\to0\) 时退回分子扩散；形式上 \(D\to0^+\) 使 \(D_{\rm eff}\to\infty\)，但混合时间 \(a^2/D\to\infty\)，先决条件 \(t\gg a^2/D\) 失效，短时仍是剪切弹道展宽。
